\backhead{FURTHER EXERCISES}{khred}
\begin{exlist}
\item For the sandy loam S05 ($\alpha=0.075$, $n=1.89$) and $\kappa=3.47\times10^{-10}$: (a) find the exponent $E=2n+\ell(n-1)$ with $\ell=0.5$; (b) find the reach $h_c$; (c) find the largest safe trigger $h_c/2$; (d) does a 600 cm trigger fire?\exans{(a) 4.225; (b) 1618 cm; (c) 809 cm; (d) yes}
\item By what factor does the reach of the sand change if $\kappa$ is ten times larger? And of the silty clay loam? Explain the difference.\exans{sand 0.69, silty clay loam 0.41 (smaller); the exponent $1/E$ is larger for small $n$}
\item An irrigation of 25 mm on a soil whose drained share follows $\eta_D=0.3(1-e^{-d/4})$ is given (a) in one event, (b) in five events of 5 mm. Find the water drained in each case.\exans{(a) 7.5 mm; (b) 5.4 mm}
\item A loam (memory fate $\phi_D=0.29$) releases 40 mm of stored water during a dry spell. How much of it drains, and how much is available to the crop and to evaporation? Repeat for the sand ($\phi_D=0.85$).\exans{loam 11.6 and 28.4 mm; sand 34 and 6 mm}
\item A sensor at 20 cm in the loam fires at 300 cm. The root zone of 0.6 m has $\theta_{fc}=0.242$ and $\theta=0.170$ at 300 cm. With $ET_c=5$ mm\,d$^{-1}$ and an efficiency of 0.85, find the interval, the net and the gross depth.\exans{8.6 d; 43 mm; 51 mm}
\end{exlist}
\backhead{ESSAY QUESTIONS}{khred}
{\small\itshape Answer each question in about one page, with equations, a sketch and at least one number from this chapter.}\par\vspace{3pt}
\begin{exlist}
\item Explain what the reach of a suction trigger is, why it exists, and why it depends on the shape of the conductivity curve rather than on the saturated conductivity.
\item A farmer uses the same 300 cm trigger on all his fields. Discuss the consequences on a sand, a loam and a clay, using the reaches of the ten soils.
\item Discuss the meaning of the memory fate of drainage and the rate-competition law. Why should one constant describe ten very different soils?
\item Explain the fate of an irrigation millimetre. Why is it important to compare an irrigated run with the rainfed run of the same soil and years?
\item Small doses drain less. Discuss the physical reason, the evidence of this study and the practical limits of very small doses.
\item Why is a fixed trigger suction not a good rule for all soils, and what rule does this chapter propose instead? Discuss its advantages and its limits.
\end{exlist}
\dykband{EXERCISES IN SWAFL}
{\small Open SWAFL (desktop or online) with the benchmark project. The chapter of the \href{https://github.com/halilunaziru73-creator/Coupled-Water-Fluxes-Pore-Scale-Dynamics-Precision-Irrigation-Mulching-Feedback-Loops/raw/main/Software/v1.1.0/SWAFL_Desktop/docs/manual_v1.1.0/SWAFL_Manual_v1.1.0.pdf}{SWAFL manual} that explains each task is given in brackets; record your values in a table and compare them with what is expected (\emph{Check}).}\par\vspace{3pt}
\begin{exlist}
\item {\color{khblue}[Manual, Chapter 6]} Open \emph{Sensor and trigger}. For the sand, set the trigger to 150, 300 and 600 cm and run \emph{All 20 years}. Record the irrigation and the yield. At which trigger does the irrigation stop?\par\noindent{\small\itshape Check: the sand is never irrigated above its reach of about 200 cm.}
\item {\color{khblue}[Manual, Chapter 6]} In the same layer, drag the purple trigger marker on the loam from 100 to 1000 cm and watch the next irrigation date and the depletion. Relate the interval to the readily available water.\par\noindent{\small\itshape Check: the interval grows with the trigger until the crop is stressed.}
\item {\color{khblue}[Manual, Chapter 6]} Open \emph{Trigger law} and read the depletion of the root zone at the trigger for three soils. Check that a single depletion fraction near 0.45 to 0.5 describes them.\par\noindent{\small\itshape Check: the triggers of all soils collapse onto one depletion fraction.}
\item {\color{khblue}[Manual, Chapter 18]} Use the search for the most efficient inputs on the drip layer of the loam with the trigger free. Which trigger does the search choose, and how does it compare with the reach?\par\noindent{\small\itshape Check: the chosen trigger lies well inside the reach.}
\end{exlist}
